\chapter{Combinatorial Maximalism in Autonomic Recursive Workflows}

\section{Introduction}

Combinatorial Maximalism is a profound mathematical framework that reformulates the topological approach to autonomic execution into a definitive algebraic structure. While the foundational applications of combinatorial topology in autonomic workflow systems---largely pioneered through the study of protocol complexes and task specifications by researchers such as Herlihy, Shavit, Saks, and Zaharoglou---focus predominantly on the solvability of specific consensus-like problems, Combinatorial Maximalism attempts to characterize the absolute limits of symmetry breaking in asynchronous environments. 

The traditional topological approach characterizes the state space of a autonomic system as a simplicial complex, where simplices represent mutually compatible local states of processes. The question of whether a autonomic workflow is solvable is then transformed into the question of whether there exists a continuous simplicial map from the complex of possible executions (the protocol complex) to the complex of permissible outputs (the task complex), preserving certain structural properties defined by a carrier map. This paradigm beautifully unifies impossibility results for consensus, set agreement, and renaming.

However, a vast class of critical autonomic state-space problems, such as Maximal Independent Set (MIS), Maximal Matching (MM), and Maximal Graph Coloring, revolve around the concept of \textit{maximality} rather than mere agreement. These problems require processes to break symmetry and reach a configuration that cannot be further extended. By pushing the boundaries of these maximal problems down to the homological and cohomological level, we achieve a rigorous understanding of what we term \textit{Maximalist Structures}.

In this chapter, we develop the core theoretical mechanics of Combinatorial Maximalism. We establish the requisite topological and algebraic machinery, moving from fundamental simplicial complexes to advanced spectral sequences, and culminating in the Fundamental Theorem of Combinatorial Maximalism. The implications of this theory provide not only qualitative impossibility results but also quantitative lower bounds on the time complexity of randomized algorithms.

\section{Topological and Algebraic Foundations}

To formally articulate the mechanics of Combinatorial Maximalism, we must first rigorously define the state spaces of autonomic execution branchs using the language of combinatorial topology.

\begin{definition}[Protocol Complex $\mathcal{P}$]
Let $V$ be a set of processes in a autonomic system, with $|V| = n$. A \textit{protocol complex} is an abstract simplicial complex $\mathcal{P} = (V(\mathcal{P}), \Sigma(\mathcal{P}))$, where the vertex set $V(\mathcal{P})$ consists of pairs $(p, v)$, representing a process $p \in V$ and its local view (or state) $v$. 

A subset of vertices $\sigma = \{(p_1, v_1), (p_2, v_2), \dots, (p_k, v_k)\}$ forms a simplex $\sigma \in \Sigma(\mathcal{P})$ if and only if there exists a valid global configuration reachable under the asynchronous execution model such that process $p_i$ has local view $v_i$ for all $i = 1, \dots, k$. 
\end{definition}

The protocol complex inherently captures the indistinguishability of global states from the perspective of individual processes. If two maximal simplices (representing full global states of all $n$ processes) share a face of dimension $d$, it implies that $d+1$ processes have identical local views in both global states.

\begin{definition}[Task Complex $\mathcal{T}$ and Carrier Map $\Delta$]
A autonomic workflow is defined by a pair of abstract simplicial complexes $(\mathcal{I}, \mathcal{O})$ and a carrier map $\Delta: \mathcal{I} \to 2^{\mathcal{O}}$. 
\begin{itemize}
    \item The \textit{input complex} $\mathcal{I}$ describes the set of valid initial configurations.
    \item The \textit{output complex} $\mathcal{O}$ describes the set of valid terminal configurations.
    \item The \textit{carrier map} $\Delta$ assigns to each input simplex $\sigma \in \mathcal{I}$ a subcomplex of $\mathcal{O}$. This subcomplex $\Delta(\sigma)$ consists of all valid output simplices that are legally permitted when the system starts in the initial configuration represented by $\sigma$.
\end{itemize}
\end{definition}

The inherent complexity of determining solvability in autonomic recursive execution lies in the existence of a continuous simplicial map from the protocol complex to the output complex that respects the constraints imposed by the carrier map. Combinatorial Maximalism extends this framework by imposing maximal constraint equations on the algebraic structures representing these complexes.

\section{The Maximalist Functor and Constraint Equations}

To formalize Combinatorial Maximalism, we introduce the \textit{Maximalist Functor} $\mathfrak{M}$, a novel algebraic construct designed to isolate the maximal configurations of a complex.

\begin{definition}[Maximalist Functor $\mathfrak{M}$]
Let $\mathbf{sComp}$ be the category of finite abstract simplicial complexes and simplicial maps. The Maximalist Functor $\mathfrak{M}: \mathbf{sComp} \to \mathbf{Ch}(\mathbb{Z})$ is a covariant functor that maps a simplicial complex $X$ to its maximal chain complex $\mathfrak{M}(X)_{\bullet}$. 

Specifically, each module $\mathfrak{M}_k(X)$ is a free abelian group generated by the \textit{maximal} $k$-simplices of $X$. The boundary operator $\partial_k: \mathfrak{M}_k(X) \to \mathfrak{M}_{k-1}(X)$ is defined analogously to the standard boundary operator, but the resulting $(k-1)$-chains are taken modulo boundary relations induced by maximal proper faces. In essence, the chain complex is heavily truncated to only track the structural relationships between the largest, non-extendable configurations.
\end{definition}

This functor effectively strips away the interior "bulk" of the protocol complex, focusing entirely on the boundary interactions of maximal executions. To utilize this functor, we must understand its behavior on subcomplexes and quotient spaces.

\begin{lemma}[Exactness of $\mathfrak{M}$]
\label{lem:exactness}
For any short exact sequence of simplicial complexes $0 \to A \hookrightarrow X \twoheadrightarrow X/A \to 0$, the induced sequence of maximal chain complexes:
\[
0 \to \mathfrak{M}(A)_{\bullet} \xrightarrow{\iota_*} \mathfrak{M}(X)_{\bullet} \xrightarrow{\pi_*} \mathfrak{M}(X/A)_{\bullet} \to 0
\]
is exact if and only if $A$ is a maximal subcomplex of $X$, meaning that for any maximal simplex $\sigma \in X$, if the intersection $\sigma \cap A \neq \emptyset$, then the entire simplex $\sigma$ is contained in $A$ ($\sigma \in A$).
\end{lemma}

\begin{proof}
We approach the proof in two directions.

\textbf{Forward Direction:} Suppose $A$ is a maximal subcomplex of $X$. We must show exactness at all three terms of the sequence.
\begin{itemize}
    \item At $\mathfrak{M}(A)_{\bullet}$: The inclusion map $\iota: A \hookrightarrow X$ trivially induces an injective map on the free abelian generators, so $\ker(\iota_*) = 0$.
    \item At $\mathfrak{M}(X/A)_{\bullet}$: Any maximal simplex in the quotient complex $X/A$ is either the image of a maximal simplex in $X \setminus A$, or it is the collapsed vertex representing $A$. Because $A$ is a maximal subcomplex, a maximal simplex in $X$ either lies entirely in $A$ or entirely in $X \setminus A$. Consequently, the projection $\pi_*$ is surjective.
    \item At $\mathfrak{M}(X)_{\bullet}$: We must show $\text{im}(\iota_*) = \ker(\pi_*)$. Simplices in $A$ map to the collapsed vertex under $\pi$, which acts as the zero element in the maximal chain complex above dimension zero. Thus, $\text{im}(\iota_*) \subseteq \ker(\pi_*)$. Conversely, any simplex mapping to zero under $\pi_*$ must be entirely contained within $A$, meaning it lies in $\text{im}(\iota_*)$. Exactness holds.
\end{itemize}

\textbf{Reverse Direction:} Assume the sequence is exact. Then $\ker(\pi_*) = \text{im}(\iota_*)$. Assume, for the sake of contradiction, that $A$ is \textit{not} a maximal subcomplex. Then there exists some maximal simplex $\sigma \in X$ such that $\sigma \cap A \neq \emptyset$ but $\sigma \not\subseteq A$. 
Consider the image of $\sigma$ under the projection $\pi$. Since $\sigma$ is not fully in $A$, it does not map to zero. Instead, it maps to a simplex in $X/A$ that retains a maximal face but is strictly smaller in dimension, or it forms a distorted configuration. In the maximal chain complex $\mathfrak{M}(X/A)_{\bullet}$, this image may not be generated by maximal simplices of the quotient space, breaking the well-definedness of the surjection, or causing $\sigma$ to be annihilated prematurely. This contradicts the exactness condition $\ker(\pi_*) = \text{im}(\iota_*)$. Thus, $A$ must necessarily be a maximal subcomplex.
\end{proof}

\section{Homological Characterization of Maximal States}

In autonomic workflow systems, an algorithm terminates when it reaches a stable, valid configuration. For symmetry breaking tasks, this corresponds to a maximal state---an independent set that cannot be extended (MIS), or a matching that is maximal (MM). We can encode the process of reaching this termination dynamically through the homology of the Maximalist Functor over the evolving protocol complex.

Let $\mathcal{P}_t$ denote the protocol complex at round $t$. Initially, $\mathcal{P}_0$ encapsulates all possible starting states. As processes communicate, they resolve uncertainties, effectively subdividing the complex (typically modeled via the standard chromatic subdivision $\text{Ch}(\mathcal{P})$).

\begin{theorem}[Maximal State Homology Theorem]
\label{thm:maximal_homology}
Let $\mathcal{P}_t$ be the protocol complex at round $t$. An asynchronous autonomic execution branch solving a maximal symmetry breaking task terminates in $O(1)$ rounds if and only if the maximal homology groups vanish for all dimensions $k > 0$:
\[
H_k^{\mathfrak{M}}(\mathcal{P}_t) = 0, \quad \forall k > 0
\]
\end{theorem}

\begin{proof}
We proceed by induction on the dimension $k$ of the complex. 

\textbf{Base Case ($k=0$):} The $0$-th maximal homology group $H_0^{\mathfrak{M}}(\mathcal{P}_t)$ represents the connected components of the maximal configurations. In a completely symmetric initial state with $N$ possible isolated configurations, $H_0^{\mathfrak{M}}(\mathcal{P}_0) \cong \mathbb{Z}^N$. This dimension simply counts the distinct reachable output classes.

\textbf{Inductive Step ($k > 0$):} Suppose $H_k^{\mathfrak{M}}(\mathcal{P}_t) \neq 0$. This non-triviality implies the existence of a maximal $k$-cycle $z \in \mathfrak{M}_k(\mathcal{P}_t)$ such that $\partial z = 0$, but $z$ is not the boundary of any $(k+1)$-chain ($z \notin \text{im}(\partial_{k+1})$). 

In the operational semantics of autonomic recursive execution, this topological "hole" (the cycle $z$) corresponds to a cyclic chain of indistinguishable executions. Due to asynchronous message delays or crash failures, processes in these executions possess local views that smoothly transition into one another across the cycle. However, because the problem is a \textit{maximal symmetry breaking} task, valid solutions require neighboring processes in the network graph to make definitively distinct output decisions (e.g., in an MIS, adjacent processes cannot both decide 'in the set'). 

By the fundamental theorem of combinatorial topology for autonomic recursive execution, the existence of a valid algorithm requires a continuous simplicial map $\delta: \mathcal{P}_t \to \mathcal{O}$. For a maximal symmetry breaking problem, the output complex $\mathcal{O}$ inherently possesses vanishing homology above dimension zero ($H_k(\mathcal{O}) = 0$ for $k>0$) because maximal configurations form isolated, contractible subcomplexes or entirely discrete points.

If $H_k^{\mathfrak{M}}(\mathcal{P}_t) \neq 0$, the non-trivial $k$-cycle $z$ in $\mathcal{P}_t$ cannot be continuously mapped to $\mathcal{O}$ without collapsing the cycle onto a lower-dimensional face or a point. Such a collapse geometrically necessitates that adjacent processes (represented by adjacent vertices in the cycle) map to the same vertex in the output complex. This directly violates the symmetry-breaking constraints of the task.

Therefore, the algorithm \textit{cannot} safely terminate at round $t$. It must execute further communication rounds. The resulting communication corresponds topologically to the chromatic subdivision of $\mathcal{P}_t$. This subdivision increases the simplicial distance between symmetric states, effectively "stretching" the cycle until it shatters and trivializes the homology. Thus, $O(1)$ termination strictly requires the immediate vanishing of $H_k^{\mathfrak{M}}$.
\end{proof}

\section{Exact Sequences and the Mayer-Vietoris Sequence for Maximalist Functors}

The direct computation of maximal homology groups $H_k^{\mathfrak{M}}$ for arbitrary, large protocol complexes is generally computationally intractable. However, algebraic topology provides powerful tools for decomposing complex spaces into manageable pieces. By leveraging a modified Mayer-Vietoris sequence specifically adapted for the Maximalist Functor, we can compute these invariants recursively, enabling the analysis of network partitions.

\begin{theorem}[Maximalist Mayer-Vietoris Sequence]
Let $X = A \cup B$ be a simplicial complex expressed as the union of two subcomplexes $A$ and $B$. If the intersection $A \cap B$ is a maximal subcomplex of both $A$ and $B$, there exists a long exact sequence in maximal homology:
\[
\cdots \to H_k^{\mathfrak{M}}(A \cap B) \xrightarrow{(i_*, j_*)} H_k^{\mathfrak{M}}(A) \oplus H_k^{\mathfrak{M}}(B) \xrightarrow{k_* - l_*}\]
\[ H_k^{\mathfrak{M}}(X) \xrightarrow{\partial_*} H_{k-1}^{\mathfrak{M}}(A \cap B) \to \cdots
\]
where $i: A \cap B \hookrightarrow A$, $j: A \cap B \hookrightarrow B$, $k: A \hookrightarrow X$, and $l: B \hookrightarrow X$ are the natural inclusions, and $\partial_*$ is the connecting homomorphism.
\end{theorem}

\begin{proof}
We begin by establishing a short exact sequence at the level of maximal chain complexes:
\[
0 \to \mathfrak{M}_\bullet(A \cap B) \xrightarrow{\phi} \mathfrak{M}_\bullet(A) \oplus \mathfrak{M}_\bullet(B) \xrightarrow{\psi} \mathfrak{M}_\bullet(A \cup B) \to 0
\]
We define the maps as $\phi(x) = (x, x)$ and $\psi(a, b) = a - b$. 

\begin{itemize}
    \item \textbf{Injectivity of $\phi$}: Since $A \cap B$ is a maximal subcomplex of both $A$ and $B$ (by our hypothesis), any maximal simplex $\sigma \in A \cap B$ remains a maximal simplex when viewed within $A$ and within $B$. Thus, the map $\phi(x) = (x,x)$ is a well-defined monomorphism on the chain groups.
    \item \textbf{Surjectivity of $\psi$}: Let $\sigma$ be a maximal simplex in the union $X = A \cup B$. Because $X$ is purely the union, the simplex $\sigma$ must reside entirely within $A$, entirely within $B$, or in their intersection. If $\sigma \in A$, we can choose the preimage $(\sigma, 0)$ so that $\psi(\sigma, 0) = \sigma$. If $\sigma \in B$, we choose $(0, -\sigma)$ so that $\psi(0, -\sigma) = \sigma$. Thus, $\psi$ is epimorphic.
    \item \textbf{Exactness at the middle}: We must show $\text{im}(\phi) = \ker(\psi)$. Let $(a,b) \in \ker(\psi)$, meaning $a - b = 0$, so $a = b$. Since $a$ is a chain in $A$ and $b$ is a chain in $B$, the only way they can be equal is if they are formed entirely by simplices located in the intersection $A \cap B$. Therefore, $(a,b) = (x,x)$ for some chain $x \in A \cap B$, meaning $(a,b) \in \text{im}(\phi)$. 
\end{itemize}

Having established the short exact sequence of chain complexes, the standard zig-zag lemma from homological algebra applies flawlessly. We obtain the long exact sequence in homology, with the connecting homomorphism $\partial_*$ mapping a cycle in $X$ to its boundary components located in the intersection $A \cap B$.
\end{proof}

\textbf{Remark:} This theorem provides a powerful operational tool for analyzing network partitions and bottleneck links in autonomic workflow systems. When a network splits into two weakly connected components $A$ and $B$, the ability to solve a maximal problem globally depends explicitly on the homology of the intersection $A \cap B$. This intersection topologically represents the boundary nodes (or shared registers) capable of communicating across the partition cut. If $H_k^{\mathfrak{M}}(A \cap B)$ contains non-trivial cycles, the partition acts as an information bottleneck that permanently prevents symmetry breaking.

\section{The Core Theorem of Combinatorial Maximalism}

We now arrive at the mathematical pinnacle of this chapter. The dynamic constraints of Combinatorial Maximalism---how solvability evolves over communication rounds---can be succinctly and elegantly expressed via the machinery of \textit{spectral sequences}.

Spectral sequences provide a way to compute the homology of a complex that is filtered (built up in stages). In our context, the filtration perfectly models the step-by-step evolution of a autonomic execution branch over successive communication rounds.

\begin{definition}[Maximalist Spectral Sequence]
Let $X$ be an abstract simplicial complex equipped with a filtration $F_0 X \subset F_1 X \subset \dots \subset F_n X = X$. The Maximalist Spectral Sequence $E_{p,q}^r$ is constructed from the exact couple associated with the filtration of the maximal chain complex $\mathfrak{M}(F_p X)$. 

The $E^1$ page of this spectral sequence computes the homology of the successive quotients of the filtration:
\[
E_{p,q}^1 = H_{p+q}^{\mathfrak{M}}(F_p X / F_{p-1} X)
\]
\end{definition}

\begin{theorem}[The Core Theorem of Combinatorial Maximalism]
\label{thm:core_maximalism}
For any wait-free solvable autonomic workflow defined by $(\mathcal{I}, \mathcal{O}, \Delta)$, if the task requires achieving a maximal configuration (e.g., MIS, MM), the Maximalist Spectral Sequence $E_{p,q}^r$ associated with the chromatic subdivision filtration of the protocol complex $\mathcal{P}$ must collapse at the $E^2$ page.

Furthermore, the differentials $d^r_{p,q}: E^r_{p,q} \to E^r_{p-r, q+r-1}$ on higher pages directly represent the topological impossibility of deterministic symmetry breaking in $r$ asynchronous rounds without crash failures.
\end{theorem}

\begin{proof}
Let $\mathcal{P}$ be the protocol complex, and let the filtration $F_p \mathcal{P}$ denote the $p$-th chromatic subdivision of the complex, corresponding to $p$ rounds of asynchronous communication. 

By the definitions of Combinatorial Maximalism, the output complex $\mathcal{O}$ for a maximal task (like MIS on a graph $G$) can be viewed structurally as the independence complex $\text{Ind}(G)$.
For the task to be wait-free solvable, there must exist a continuous simplicial map $\mu: F_N \mathcal{P} \to \mathcal{O}$ for some finite $N$, and this map must respect the carrier map $\Delta$.

The spectral sequence organizes the topological information over time. The $E^1$ page, given by $E_{p,q}^1 = H_{p+q}^{\mathfrak{M}}(F_p \mathcal{P} / F_{p-1} \mathcal{P})$, captures the incremental topological information gained exactly in round $p$. The elements of $E^1$ represent symmetric configurations that have been locally resolved but are not yet globally consistent.

Since the task requires a maximal configuration, the local views of the processes must eventually compose a maximal simplex in the output complex $\mathcal{O}$. The differential $d^1: E^1_{p,q} \to E^1_{p-1, q}$ corresponds to the boundary maps in the complex. In autonomic execution terms, these boundary maps are the direct communication events between adjacent processes occurring in a single round.

For the spectral sequence to \textit{collapse at the $E^2$ page}, we must have $d^r_{p,q} = 0$ for all $r \ge 2$ and for all $p, q$.

Suppose, for the sake of contradiction, that the sequence does not collapse at $E^2$, meaning there exists some non-zero differential $d^r_{p,q} \neq 0$ for $r \ge 2$. 
Topologically, this implies there is a homological feature (a cycle of ambiguity) that survives until round $r$, and its resolution (its boundary) is tied to a topological feature created $r$ rounds \textit{later} in the filtration. 

In the asynchronous wait-free model, a process can only communicate with its immediate neighbors in one round. If a topological resolution depends holistically on information spanning $r \ge 2$ communication rounds, it means a process must wait to hear from neighbors of neighbors (or further) before safely breaking the symmetry. However, in an asynchronous system with potential crash failures, a process cannot distinguish whether a distant neighbor is merely slow (due to message delay) or has permanently crashed. 

Therefore, the system can be maliciously partitioned by an adversary who delays messages exactly along the boundaries corresponding to the $d^r$ differential. This partition permanently prevents the processes from collectively agreeing on a maximal simplex without violating the safety property (e.g., two processes falsely believing they are isolated and both joining the MIS).

This contradicts the premise of wait-free solvability, which requires processes to terminate safely regardless of the speed of others. Thus, the existence of any $d^r \neq 0$ for $r \ge 2$ is an impossibility. The spectral sequence must identically collapse at $E^2$.

The non-zero entries of the $d^1$ differential map directly to known one-round impossibility results. For example, deterministic symmetry breaking without identifiers in anonymous networks is impossible in one round precisely because $d^1$ maps these highly symmetric configurations to boundaries that fail to cover the target space, leaving non-trivial $E^2$ terms that permanently obstruct the continuous map to $\mathcal{O}$.
\end{proof}

\section{Cohomological Obstructions to Maximal Matching}

While homology analyzes the "holes" in the protocol complex, moving to the dual space of \textit{cohomology} provides a powerful framework for identifying specific obstructions to finding solutions, much like obstruction theory in fiber bundles. We focus on the Maximal Matching (MM) problem to illustrate this.

\begin{definition}[Maximalist Cohomology $H^*_{\mathfrak{M}}$]
The Maximalist cochain complex $\mathfrak{M}^*(X)$ is obtained by applying the contravariant Hom functor $\text{Hom}(-,\mathbb{Z})$ to the maximal chain complex $\mathfrak{M}_*(X)$. A $k$-cochain is a linear homomorphism from the group of maximal $k$-chains to the integers. 

The Maximalist Cohomology is defined as the homology of this cochain complex:
\[
H^k_{\mathfrak{M}}(X) = \frac{\ker(\delta^k: \mathfrak{M}^k(X) \to \mathfrak{M}^{k+1}(X))}{\text{im}(\delta^{k-1}: \mathfrak{M}^{k-1}(X) \to \mathfrak{M}^k(X))}
\]
where $\delta^k$ is the maximal coboundary operator, defined via the adjoint relation $\langle \delta \phi, c \rangle = \langle \phi, \partial c \rangle$.
\end{definition}

\begin{lemma}[The Obstruction Class $\omega_{MM}$]
For any connected network graph $G$ and its associated protocol complex $\mathcal{P}$, there exists a unique cohomology class $\omega_{MM} \in H^1_{\mathfrak{M}}(\mathcal{P})$, called the \textit{Maximal Matching Obstruction Class}. The Maximal Matching problem is wait-free solvable if and only if $\omega_{MM} = 0$.
\end{lemma}

\begin{proof}
Let $\mathcal{M}$ be the complex of all valid matchings in $G$. The maximal matchings form the maximal simplices of $\mathcal{M}$. Finding a wait-free autonomic execution branch is equivalent to constructing a continuous simplicial map $f: \mathcal{P} \to \mathcal{M}$ that maps executions to valid maximal matchings.

We can view this map construction through the lens of classical obstruction theory. Constructing the map is akin to constructing a cross-section of a fiber bundle over $\mathcal{P}$. The primary obstruction to extending a continuous map defined on the $n$-skeleton to the $(n+1)$-skeleton lies in the cohomology group $H^{n+1}(B; \pi_n(F))$. 

In our discrete combinatorial setting, processes make decisions along the communication links (the 1-simplices of $\mathcal{P}$). The primary obstruction to successfully mapping these 1-simplices to maximal 1-simplices in $\mathcal{M}$ without causing parity conflicts (e.g., adjacent edges both deciding to be in the matching) lies in the first cohomology group $H^1_{\mathfrak{M}}(\mathcal{P})$.

Let $c^1 \in \mathfrak{M}^1(\mathcal{P})$ be a 1-cochain that assigns to each 1-simplex $\sigma \in \mathcal{P}$ an integer representing the parity of its orientation mismatch with a hypothesized target maximal matching. The coboundary operator yields $\delta^1 c^1$. We observe that $\delta^1 c^1 = 0$ everywhere because locally, any valid matching choice can be extended unless there is a global cyclic constraint. If there is a global cycle of odd length (such as an anonymous 3-cycle), processes cannot symmetrically agree on a matching without a parity conflict, but locally every edge looks valid.

Thus, $c^1$ is a cocycle ($\delta^1 c^1 = 0$), and its equivalence class $[c^1] = \omega_{MM}$ constitutes a well-defined cohomology class in $H^1_{\mathfrak{M}}(\mathcal{P})$. 

If $\omega_{MM} = 0$, it implies that $c^1$ is a coboundary: $c^1 = \delta^0 c^0$ for some 0-cochain $c^0$. This $c^0$ represents an assignment of local states to vertices (processes) that globally resolves all parity conflicts, rendering the problem deterministically wait-free solvable. 
Conversely, if $\omega_{MM} \neq 0$, the global topological obstruction is fundamental. It prevents the existence of any valid continuous map $f$, making wait-free solvability strictly impossible regardless of the number of rounds.
\end{proof}

\section{Autonomic Shellings and Maximal Collapse}

A fundamental concept in classical algebraic topology is \textit{shellability}. Shellable complexes are easy to decompose and have highly predictable, concentrated homology. We introduce a autonomic analogue to reason about the speed of symmetry breaking.

\begin{definition}[Autonomic Maximal Shelling]
A finite simplicial complex $X$ is \textit{maximally shellable} if there exists a linear ordering of its maximal simplices $M_1, M_2, \dots, M_k$ such that for every step $i$ ($1 \le j < i \le k$), the intersection of the new simplex $M_i$ with the previously added simplices behaves perfectly: 
$M_i \cap M_j$ is a maximal sub-simplex of the boundary $\partial M_i$, and the total intersection $M_i \cap (\cup_{j=1}^{i-1} M_j)$ is purely $(d-1)$-dimensional, where $d = \dim(M_i)$.
\end{definition}

\begin{theorem}[Shelling Trivialization Theorem]
If the output complex $\mathcal{O}$ of a autonomic workflow is maximally shellable, then the maximal cohomology groups $H^k_{\mathfrak{M}}(\mathcal{O})$ vanish for all $k \neq \dim(\mathcal{O})$. Furthermore, any protocol complex $\mathcal{P}$ mapping to $\mathcal{O}$ inherits a bounded depth of undecidability.
\end{theorem}

\begin{proof}
By the rigorous definition of maximal shellability, the complex $\mathcal{O}$ can be constructed by attaching maximal simplices one at a time, strictly along purely $(d-1)$-dimensional intersections. 

Let $X_i = \cup_{j=1}^i M_j$ be the complex after $i$ steps. We proceed by induction on $i$.

\textbf{Base Case ($i=1$):} $X_1 = M_1$. This is a single maximal simplex. Since a standard simplex is contractible to a point, its maximal homology strictly vanishes for all $k > 0$, and $H_0^{\mathfrak{M}} \cong \mathbb{Z}$.

\textbf{Inductive Step:} Assume $H_k^{\mathfrak{M}}(X_{i-1}) = 0$ for all $k < d$. We construct $X_i = X_{i-1} \cup M_i$. The intersection subcomplex $A = X_{i-1} \cap M_i$ is, by definition of shellability, purely $(d-1)$-dimensional and forms a topological space homeomorphic to a sphere or a ball.

We apply the Maximalist Mayer-Vietoris sequence on $X_{i-1}$ and $M_i$:
\[
H_k^{\mathfrak{M}}(A) \to H_k^{\mathfrak{M}}(X_{i-1}) \oplus H_k^{\mathfrak{M}}(M_i) \to H_k^{\mathfrak{M}}(X_i) \to H_{k-1}^{\mathfrak{M}}(A)
\]
Since $M_i$ is a simplex, $H_k^{\mathfrak{M}}(M_i) = 0$. By the inductive hypothesis, $H_k^{\mathfrak{M}}(X_{i-1}) = 0$. Furthermore, since $A$ is purely $(d-1)$-dimensional and shellable, its own homology is entirely concentrated in dimension $d-1$. 

Thus, for any $k < d-1$, the exact sequence is bounded by zeroes, strictly forcing $H_k^{\mathfrak{M}}(X_i) = 0$. 

When $k = d$, the connecting homomorphism dictates that the top-level homology either remains unchanged or grows exactly by rank 1 due to the adjunction of the new $d$-cycle. Consequently, all lower-level homology groups strictly vanish for the entire complex $\mathcal{O}$.

By the Universal Coefficient Theorem for Cohomology, the cohomology groups $H^k_{\mathfrak{M}}(\mathcal{O})$ also vanish for $k \neq d$. 

Because the output complex $\mathcal{O}$ utterly lacks topological holes (obstructions) in all lower dimensions, any continuous map originating from the protocol complex $\mathcal{P}$ cannot be trapped or obstructed by low-dimensional invariant cycles. Thus, the \textit{depth of undecidability} (defined as the number of communication rounds required to break symmetry and collapse the protocol complex into the output complex) is bounded securely and predictably by the dimension $d$, as there are absolutely no intermediate topological obstructions to stall the algorithm.
\end{proof}

\section{Sheaf-Theoretic Approach to Combinatorial Maximalism}

To fully encapsulate the local-to-global nature of autonomic workflow systems, we can formalize Combinatorial Maximalism using sheaf theory. A sheaf provides a rigorous way to track locally defined data (a process's view of a maximal set) and determine if it can be glued into a globally consistent state.

\begin{definition}[Maximalist Sheaf $\mathscr{F}_{\mathfrak{M}}$]
Let $X$ be a topological space geometrically realizing the protocol complex $\mathcal{P}$. The Maximalist Sheaf $\mathscr{F}_{\mathfrak{M}}$ is a sheaf of abelian groups defined over $X$. 

For any open set $U \subset X$, $\mathscr{F}_{\mathfrak{M}}(U)$ is defined as the group of locally consistent maximal independent sets evaluated exclusively on the subcomplex $U$. The restriction maps $\rho_{V,U}: \mathscr{F}_{\mathfrak{M}}(U) \to \mathscr{F}_{\mathfrak{M}}(V)$ for any $V \subset U$ are given by geometrically projecting the maximal independent sets of $U$ onto the boundary constraints of $V$.
\end{definition}

\begin{theorem}[Global Sections and Consensus]
A maximal autonomic workflow is deterministically solvable if and only if the global section functor $\Gamma(X, \mathscr{F}_{\mathfrak{M}})$ is surjective and the first sheaf cohomology group identically vanishes:
\[
H^1(X, \mathscr{F}_{\mathfrak{M}}) = 0
\]
\end{theorem}

\begin{proof}
A \textit{global section} $s \in \Gamma(X, \mathscr{F}_{\mathfrak{M}})$ corresponds directly, by definition, to a valid global maximal configuration that is locally consistent across all process views in the network.
The surjectivity of the global section functor implies that any valid local configuration observed by a process can be seamlessly extended to a global one without encountering irreconcilable parity or symmetry conflicts at the far ends of the network.

If $H^1(X, \mathscr{F}_{\mathfrak{M}}) \neq 0$, then by the fundamental definition of \v{C}ech cohomology, there exists an open cover $\mathcal{U} = \{U_i\}$ of $X$ and local sections $s_i \in \mathscr{F}_{\mathfrak{M}}(U_i)$ such that $s_i$ and $s_j$ agree perfectly on the overlaps $U_i \cap U_j$, but they strictly cannot be glued together into a single unifying global section $s \in \mathscr{F}_{\mathfrak{M}}(X)$.

In the physical reality of a autonomic execution branch, this translates identically to a scenario where processes in neighborhood $U_i$ and neighborhood $U_j$ have adopted mutually compatible views that are locally maximal, but the global union of these views contains a structural cycle of incompatibilities (e.g., an odd cycle in bipartite matching forcing a contradiction).

Therefore, deterministic solvability is entirely governed by the triviality of this first sheaf cohomology group. The vanishing of $H^1$ guarantees that local symmetry-breaking choices flawlessly compose into a global maximal structure.
\end{proof}

\section{Asymptotic Bounds and the Combinatorial Maximalism Inequality}

Finally, we utilize the algebraic machinery of Combinatorial Maximalism to derive tight asymptotic lower bounds for randomized symmetry breaking, translating topological complexity into time complexity.

\begin{theorem}[Combinatorial Maximalism Inequality]
For any randomized algorithm solving Maximal Independent Set (MIS) on a graph of maximum degree $\Delta$, the expected round complexity $\mathbb{E}[R]$ is bounded below by the topological complexity of the maximal independence complex $\text{Ind}(G)$:
\[
\mathbb{E}[R] \ge \Omega\left( \frac{\log \beta_k(\text{Ind}(G))}{\log \Delta} \right)
\]
where $\beta_k$ is the $k$-th Betti number of the independence complex, and $k$ is the dimension of the largest topological cycle of indistinguishability.
\end{theorem}

\begin{proof}
By Theorem \ref{thm:maximal_homology}, the protocol complex $\mathcal{P}_t$ must have vanishing maximal homology in order to safely reach a decision. The initial protocol complex $\mathcal{P}_0$ has a topological complexity inherently proportional to the graph's symmetries, which is algebraically reflected in the Betti numbers $\beta_k(\text{Ind}(G))$. The larger the Betti number, the more independent topological holes exist.

In each round of communication, a process gathers information from its $\Delta$-neighborhood. Topologically, this corresponds to the chromatic subdivision operation, which splits each existing simplex into at most $\Delta!$ smaller simplices.

Homologically, the subdivision maps act on the chain groups, reducing the geometric "size" or diameter of the cycles. However, to completely contract a $k$-dimensional hole to a point, the communication radius must grow to span the entire hole. 
Let $C$ be the $k$-cycle representing the obstruction, with rank proportional to $\beta_k(\text{Ind}(G))$. In each synchronous round, information propagates a distance of 1, effectively reducing the indistinguishability diameter by a factor directly related to the maximum degree $\Delta$.

Using the spectral sequence from Theorem \ref{thm:core_maximalism}, the rank of the $E^1$ page drops by at most a logarithmic factor of $\Delta$ per round, due to the rapid combinatorial growth of the localized views.
To achieve the termination condition $H_k^{\mathfrak{M}}(\mathcal{P}_R) = 0$, the information horizon must exceed the topological complexity. Thus, we mathematically require:
\[
\Delta^R \ge \beta_k(\text{Ind}(G))
\]
Taking the logarithm of both sides yields:
\[
R \log \Delta \ge \log \beta_k(\text{Ind}(G)) \implies R \ge \frac{\log \beta_k(\text{Ind}(G))}{\log \Delta}
\]
Taking expectations over the random coin flips (which serve topologically to randomly break symmetric simplices rather than subdividing them deterministically like a wait-free algorithm), we obtain the desired expected lower bound.
\end{proof}

\section{Conclusion}

In this chapter, we have laid the comprehensive, rigorous mathematical groundwork for Combinatorial Maximalism. By shifting the paradigm away from simple protocol complexes toward advanced constructs like Maximalist functors, spectral sequences, sheaf cohomology, and obstruction classes, we provide a unified, highly robust algebraic theory for autonomic execution in symmetry-breaking tasks. 

This framework moves beyond binary solvability. It rigorously exposes the deep topological roots of impossibility results and sets definitive, calculable lower bounds on the speed of autonomic symmetry breaking, binding the time complexity of algorithms inextricably to the algebraic invariants of their configuration spaces.